\documentclass[12pt]{article}
\usepackage{amsmath, amssymb, array, booktabs, caption}
\usepackage{geometry}% 1-inch margins
\usepackage{enumitem} % Include this in the preamble of your LaTeX document
\geometry{letterpaper, margin=1in}% 1-inch margins
\usepackage{multicol}


\title{MATH 374 (Discrete Structures)\\Chapter 1.3 \& 1.4 Homework}
\author{Jacob Vaught}
\date{09-20-2023}


\begin{document}
\maketitle
\newpage
\section*{Problem 1.3.1}
What is the truth value of the following wffs in the interpretation where the domain consists of the integers, \(O(x)\) is ``\(x\) is odd", \(L(x)\) is ``\(x < 10\)", and \(G(x)\) is ``\(x > 9\)"?
\begin{itemize}
    \item[a.] \((\exists x)O(x)\)
    \item[b.] \((\forall x)[L(x) \implies O(x)]\)
    \item[c.] \((\exists x)[L(x) \land G(x)]\)
    \item[d.] \((\forall x)[L(x) \lor G(x)]\)
\end{itemize}

\subsection*{Solutions}

\subsubsection*{a. \((\exists x)O(x)\)}

The statement asserts that there exists at least one integer \( x \) such that \( x \) is odd. This is clearly true, as there are many integers that are odd, for example, \( x = 1 \), \( x = -3 \), etc.

\textbf{Truth Value:} True

\subsubsection*{b. \((\forall x)[L(x) \implies O(x)]\)}

The statement claims that for every integer \( x \), if \( x < 10 \), then \( x \) must be odd. This is not true because integers less than 10 include both odd and even numbers. For example, \( x = 2 \) is less than 10 but is not odd.

\textbf{Truth Value:} False

\subsubsection*{c. \((\exists x)[L(x) \land G(x)]\)}

The statement says there exists an integer \( x \) such that \( x < 10 \) and \( x > 9 \) simultaneously. This is impossible because no integer can be both less than 10 and greater than 9 at the same time.

\textbf{Truth Value:} False

\subsubsection*{d. \((\forall x)[L(x) \lor G(x)]\)}

The statement asserts that for every integer \( x \), \( x \) is either less than 10 or greater than 9. This statement is true because all integers can be less than 10, or larger than 9.

\textbf{Truth Value:} True


\newpage
\section*{Problem 1.3.2}
What is the truth value of each of the following wffs in the interpretation where the domain consists of the integers?

\begin{multicols}{2}
\begin{itemize}
    \item[a.] \((\forall x)(\exists y)(x + y = x)\)
    \item[b.] \((\exists y)(\forall x)(x + y = x)\)
    \item[c.] \((\forall x)(\exists y)(x + y = 0)\)
    \item[d.] \((\exists y)(\forall x)(x + y = 0)\)
    \item[e.] \((\forall x)(\forall y)(x < y \lor y < x)\)
    \item[f.] \((\forall x)[x < 0 \implies (\exists y)(y > 0 \land x + y = 0)]\)
    \item[g.] \((\exists y)(\exists y)(x^2 = y)\)
    \item[h.] \((\forall x)(x^2 > 0)\)
\end{itemize}
\end{multicols}
\subsection*{Solutions}

\subsubsection*{a. \((\forall x)(\exists y)(x + y = x)\)}
For all integers \( x \), there exists an integer \( y \) such that \( x + y = x \). This is true because you can always pick \( y = 0 \) for any \( x \).

\textbf{Truth Value:} True

\subsubsection*{b. \((\exists y)(\forall x)(x + y = x)\)}
There exists an integer \( y \) such that for all integers \( x \), \( x + y = x \). This is true because you can pick \( y = 0 \).

\textbf{Truth Value:} True

\subsubsection*{c. \((\forall x)(\exists y)(x + y = 0)\)}
For every integer \( x \), there exists an integer \( y \) such that \( x + y = 0 \). This is true because for any \( x \), you can pick \( y = -x \).

\textbf{Truth Value:} True

\subsubsection*{d. \((\exists y)(\forall x)(x + y = 0)\)}
There exists an integer \( y \) such that for all integers \( x \), \( x + y = 0 \). This is false because there is no single integer \( y \) that would make \( x + y = 0 \) true for all integers \( x \).

\textbf{Truth Value:} False

\subsubsection*{e. \((\forall x)(\forall y)(x < y \lor y < x)\)}
For all integers \( x \) and \( y \), \( x < y \) or \( y < x \). This is false because it doesn't account for the case where \( x = y \).

\textbf{Truth Value:} False

\subsubsection*{f. \((\forall x)[x < 0 \implies (\exists y)(y > 0 \land x + y = 0)]\)}
For all integers \( x \) that are less than 0, there exists an integer \( y \) such that \( y > 0 \) and \( x + y = 0 \). This is true because for any \( x < 0 \), you can pick \( y = -x \), which is greater than 0.

\textbf{Truth Value:} True

\subsubsection*{g. \((\exists x)(\exists y)(x^2 = y)\)}
There exist integers \( x \) and \( y \) such that \( x^2 = y \). The statement is true, as for any given integer \( x \), the value of \( x^2 \) will also be an integer, fulfilling the condition for some \( y \).

\textbf{Truth Value:} True


\subsubsection*{h. \((\forall x)(x^2 > 0)\)}
For all integers \( x \), \( x^2 > 0 \). This is false because it doesn't account for \( x = 0 \), where \( x^2 = 0 \).

\textbf{Truth Value:} False
\newpage
\section*{Problem 1.3.5}
For each wff, find an interpretation in which it is true and one in which it is false.
\begin{itemize}
    \item[a.] \((\forall x)([A(x) \lor B(x)] \land \lnot [A(x) \land B(x)])\)
    \item[b.] \((\forall x)(\forall y)[P(x, y) \implies P(y, x)]\)
    \item[c.] \((\forall x)[P(x) \implies (\exists y)Q(x, y)]\)
    \item[d.] \((\exists x)[A(x) \land (\exists y)B(x, y)]\)
    \item[e.] \([( \forall x)A(x) \implies (\forall x)B(x)] \implies (\forall x)[A(x) \implies B(x)]\)
\end{itemize}

\subsection*{Solutions}

\subsubsection*{a. \((\forall x)([A(x) \lor B(x)] \land \lnot [A(x) \land B(x)])\)}

\textbf{True Interpretation:} Let \(A(x)\) mean "x is an apple" and \(B(x)\) mean "x is an orange". In the domain of fruits, this statement holds true because any fruit is either an apple or an orange, but not both.
\newline
\textbf{False Interpretation:} Let \(A(x)\) mean "x is a citrus fruit" and \(B(x)\) mean "x is orange-colored". In this domain, the statement is false when \(x\) is an orange, because an orange is both a citrus fruit and orange-colored.


\subsubsection*{b. \((\forall x)(\forall y)[P(x, y) \implies P(y, x)]\)}

\textbf{True Interpretation:} Let \(P(x, y)\) mean "x and y are reachable from each other in an undirected graph" over the domain of nodes in the graph. In this case, the statement holds true. If node \(x\) is reachable from node \(y\) in an undirected graph, then by definition, \(y\) is also reachable from \(x\).
\newline
\textbf{False Interpretation:} Let \(P(x, y)\) mean "there is a directed edge from \(x\) to \(y\)" over the domain of nodes in a directed graph. This statement is false because having a directed edge from \(x\) to \(y\) does not necessarily imply that there is a directed edge from \(y\) to \(x\).

\subsubsection*{c. \((\forall x)[P(x) \implies (\exists y)Q(x, y)]\)}

\textbf{True Interpretation:} Let \(P(x)\) mean "x is a computer made after 2015" and \(Q(x, y)\) mean "y is an operating system compatible with \(x\)." In this interpretation, for every computer made after 2015, there exists an operating system that is compatible with it, making the statement true.
\newline
\textbf{False Interpretation:} Let \(P(x)\) mean "x is a computer made before 2015" and \(Q(x,y)\) mean "y is a computer that ran Windows 11 at release." This statement is false because a computer made before 2015 couldn't run Windows 11 before it's 2021 release date.


\subsubsection*{d. \((\exists x)[A(x) \land (\forall y)B(x, y)]\)}

\textbf{True Interpretation:}
Let \(A(x)\) denote "x is an odd number" and \(B(x, y)\) signify "x + y is an integer." In this case, the statement is true, for instance, for \(x=1\). Any \(y\) added to \(x\) would result in an integer.
\newline
\textbf{False Interpretation:}
Let \(A(x)\) mean "x is a prime number" and \(B(x, y)\) denote "x + y is divisible by 3." The statement is false because, for example, 2 is a prime number, but \(2 + 1\) is not divisible by 3, invalidating \(\forall y, B(x, y)\) in this context.


\subsubsection*{e. \([( \forall x)A(x) \implies (\forall x)B(x)] \implies (\forall x)[A(x) \implies B(x)]\)}

\textbf{True Interpretation:} Let \(A(x)\) mean "x is a functioning computer" and \(B(x)\) mean "x can run an instruction" over the domain of computers. This statement is true because if a computer is functioning, it must be able to run an instruction.
\newline
\textbf{False Interpretation:} Let \(A(x)\) mean "x is a functioning computer" and \(B(x)\) mean "x runs Windows 11" over the domain of computers. This statement is false because if a computer is functioning, it does not necessarily run Windows 11.


\newpage
\section*{Problem 1.3.16}
Using the predicate symbols shown and the appropriate quantifiers, write each English language statement as a predicate wff. (The domain is the whole world.)

\( B(x) \) is ``\( x \) is a bee''.

\( F(x) \) is ``\( x \) is a flower''.

\( L(x, y) \) is ``\( x \) loves \( y \)''.

\begin{enumerate}
    \item[a.] All bees love all flowers.
    \item[b.] Some bees love all flowers.
    \item[e.] Only bees love flowers.
    \item[g.] No bee loves only flowers.
\end{enumerate}
\subsection*{Solutions}
\begin{enumerate}
    \item[a.] All bees love all flowers.
    
    This statement can be written as:
    \[
    (\forall x)(B(x) \rightarrow (\forall y)(F(y) \rightarrow L(x, y)))
    \]
    This wff states that for every \( x \) that is a bee, for every \( y \) that is a flower, \( x \) loves \( y \).
    
    \item[b.] Some bees love all flowers.
    
    This statement can be written as:
    \[
    (\exists x)(B(x) \land (\forall y)(F(y) \rightarrow L(x, y)))
    \]
    This wff says that there exists at least one \( x \) that is a bee, for which every \( y \) that is a flower is loved by \( x \).
    
    \item[e.] Only bees love flowers.
    
    This statement can be written as:
    \[
    (\forall x)(\forall y)(L(x, y) \land F(y) \rightarrow B(x))
    \]
    This wff states that for all \( x \) and \( y \), if \( x \) loves \( y \) and \( y \) is a flower, then \( x \) must be a bee.

    \item[g.] No bees only love flowers.
    
    This statement can be written as:
    \[
    \lnot (\exists x)(B(x) \land (\forall y)(L(x, y) \rightarrow F(y)))
    \]
    
    This wff states that there is not a bee \( x \) for which every \( y \) it loves is a flower, effectively capturing the notion that "No bee loves only flowers."
\end{enumerate}

\newpage
\section*{Problem 1.3.24}
Explain why each wff is valid.
\begin{itemize}
    \item[a.] \((\forall x)(\forall y)A(x, y) \iff (\forall y)(\forall x) A(x, y)\)
    \item[c.] \((\exists x)(\forall y)P(x, y) \implies (\forall y)(\exists x)P(x, y)\)
    \item[d.] \(A(a) \implies (\exists x)A(x)\)
\end{itemize}

\subsubsection*{Solutions}

\begin{itemize}
    \item[a.] \((\forall x)(\forall y)A(x, y) \iff (\forall y)(\forall x) A(x, y)\)
    
    \textbf{Explanation:}
    
    This wff is valid because the ordering of universal quantifiers (\( \forall x \), \( \forall y \)) doesn't matter when they are applied to the same expression \( A(x, y) \). Whether you say "for all \( x \) and for all \( y \)," or "for all \( y \) and for all \( x \)," you're still talking about every possible combination of \( x \) and \( y \) satisfying \( A(x, y) \). Hence, the statement is valid by the commutative property of the universal quantifiers over the same expression.
    
    \item[c.] \((\exists x)(\forall y)P(x, y) \implies (\forall y)(\exists x)P(x, y)\)
    
    \textbf{Explanation:}
    
    This wff says that if there exists some \( x \) such that for all \( y \), \( P(x, y) \) holds, then for every \( y \), there must exist an \( x \) for which \( P(x, y) \) holds. In simpler terms, if there's one \( x \) that works for every \( y \), then it must be true that for each individual \( y \), there's at least one \( x \) that works. Therefore, the wff is valid because the initial condition implies the latter condition.
    
    \item[d.] \(A(a) \implies (\exists x)A(x)\)
    
    \textbf{Explanation:}
    
    This is a straightforward implication. It states that if \( A(a) \) holds for some specific entity \( a \), then there must exist at least one entity \( x \) for which \( A(x) \) holds. Since \( a \) is an entity and \( A(a) \) is true, \( a \) itself satisfies the condition \( (\exists x)A(x) \), making the statement valid.
    
\end{itemize}

Each of these wffs is valid due to the specific properties and logical relationships of the quantifiers and statements involved.
\newpage
\section*{Problem 1.3.25}
Give interpretations to prove that each of the following wffs is not valid:
\begin{itemize}
    \item[a.] \((\exists x)A(x) \land (\exists x)B(x) \implies (\exists x)[A(x) \land B(x)]\)
    \item[b.] \((\forall x)(\exists y)P(x, y) \implies (\exists x)(\forall y)P(x, y)\)
    \item[c.] \((\forall x)[P(x) \implies Q(x)] \implies [(\exists x)P(x) \implies (\forall x)Q(x)]\)
    \item[d.] \((\forall x)\lnot A(x) \iff \lnot (\forall x)A(x)\)
\end{itemize}

\subsubsection*{Solutions}

\begin{itemize}
    \item[a.] \((\exists x)A(x) \land (\exists x)B(x) \implies (\exists x)[A(x) \land B(x)]\)
    
    \textbf{Interpretation:}
    
    Let \(A(x)\) mean "x is an even number," and \(B(x)\) mean "x is an odd number." In this case, \( (\exists x)A(x) \) is true because even numbers exist, and \( (\exists x)B(x) \) is true because odd numbers exist. However, \( (\exists x)[A(x) \land B(x)] \) is false because no number is both even and odd. This proves the wff is not valid.

    \item[b.] \((\forall x)(\exists y)P(x, y) \implies (\exists x)(\forall y)P(x, y)\)
    
    \textbf{Interpretation:}
    
    Let \(P(x, y)\) mean "x is smaller than y." The first part \( (\forall x)(\exists y)P(x, y) \) is true because for every number \( x \), there exists a number \( y \) greater than it. However, \( (\exists x)(\forall y)P(x, y) \) is false because there is no single number \( x \) that is smaller than every number \( y \). Therefore, the wff is not valid.

    \item[c.] \((\forall x)[P(x) \implies Q(x)] \implies [(\exists x)P(x) \implies (\forall x)Q(x)]\)
    
    \textbf{Interpretation:}
    
    Let \(P(x)\) mean "x is a dog," and \(Q(x)\) mean "x is an animal." While it's true that all dogs are animals (\(\forall x)[P(x) \implies Q(x)]\), if there exists at least one dog (\(\exists x)P(x)\), it doesn't mean that every existing thing is an animal (\(\forall x)Q(x)\). Thus, the wff is not valid.
    
    \item[d.] \((\forall x)\lnot A(x) \iff \lnot (\forall x)A(x)\)
    
    \textbf{Interpretation:}
    
    Let \(A(x)\) mean "x is happy." The statement \( (\forall x)\lnot A(x) \) means that nobody is happy. The statement \( \lnot (\forall x)A(x) \) means that it's not the case that everyone is happy. However, the latter doesn't necessarily mean that nobody is happy; it could also mean that some are and some aren't. Therefore, the wff is not valid.
\end{itemize}



\newpage
\section*{Problem 1.4.6}

Justify each step in the following proof sequence of
\[
(\exists x)[P(x)\rightarrow Q(x)] \rightarrow [(\forall x)P(x)\rightarrow (\exists x)Q(x)]
\]

\begin{enumerate}
    \item \( (\exists x)[P(x)\rightarrow Q(x)] \)
    \item \( P(a)\rightarrow Q(a) \)
    \item \( (\forall x)P(x) \)
    \item \( P(a) \)
    \item \( Q(a) \)
    \item \( (\exists x)Q(x) \)
\end{enumerate}
\subsection*{Solutions}
\begin{enumerate}
    \item \( (\exists x)[P(x)\rightarrow Q(x)] \) \hfill hypothesis
    \item \( P(a)\rightarrow Q(a) \) \hfill 1, ei [Existential instantiation]
    \item \( (\forall x)P(x) \) \hfill hypothesis
    \item \( P(a) \) \hfill 3, ui [Universal instantiation]
    \item \( Q(a) \) \hfill 2, 4, mp [Modus ponens]
    \item \( (\exists x)Q(x) \) \hfill 5, eg [Existential generalization]
\end{enumerate}
\newpage
\section*{Problem 1.4.9}
\textbf{Problem Statement:}
Consider the wff \( (\forall y)(\exists x)Q(x,y) \rightarrow (\exists x)(\forall y)Q(x,y) \).

\subsection*{a. Find an interpretation to prove that this wff is not valid.}

Consider the domain of all integers and let \( Q(x, y) \) represent \( x = y + 1 \).
\newline\newline
For the left-hand side of the wff \( (\forall y)(\exists x)Q(x,y) \):
\newline
- Given any integer \( y \), there exists an integer \( x = y + 1 \) that makes \( Q(x, y) \) true.
\newline\newline
However, for the right-hand side \( (\exists x)(\forall y)Q(x,y) \):
\newline
- There is no single integer \( x \) for which \( Q(x, y) \) will hold for all \( y \) (i.e., there is no integer \( x \) such that \( x = y + 1 \) for all \( y \)).
\newline \newline
Thus, the wff is not valid under this interpretation.

\subsection*{b. Find the flaw in the following "proof" of this wff.}

\begin{enumerate}[left=0em,labelwidth=4em,align=left]
    \item[\( (\forall y)(\exists x)Q(x,y) \)] \hfill hyp
    \item[\( (\exists x)Q(x,y) \)] \hfill 1, ui [Universal instantiation]
    \item[\( Q(a, y) \)] \hfill 2, ei [Existential instantiation]
    \item[\( (\forall y)Q(a,y) \)] \hfill 3, ug [Universal generalization]
    \item[\( (\exists x)(\forall y)Q(x,y) \)] \hfill 4, eg [Existential generalization]
\end{enumerate}
The flaw lies in step 2. Universal Instantiation is used incorrectly here, where Existential Instantiation should be used.

\newpage
\section*{Exercises 1.4.10-14}

Prove that each wff is a valid argument.

\begin{enumerate}
    \item[\textbf{10.}] \( (\forall x)P(x) \rightarrow (\forall x)[P(x)\lor Q(x)] \)
    
    \item[\textbf{11.}] \( (\forall x)P(x)\land (\exists x)Q(x)\rightarrow (\exists x)[P(x)\land Q(x)] \)
    
    \item[\textbf{12.}] \( (\exists x)(\exists y)P(x,y)\rightarrow (\exists y)(\exists x)P(x,y) \)
\end{enumerate}

\section*{Solutions to Exercises 1.4.10-14}

\subsection*{Exercise 10}
Prove that the argument \( (\forall x)P(x) \rightarrow (\forall x)[P(x)\lor Q(x)] \) is valid.

\subsubsection*{Solution}
\begin{enumerate}[left=0em,labelwidth=4em,align=left]
    \item[\( (\forall x)P(x) \)] \hfill Premise
    \item[\( P(a) \)] \hfill 1, UI [Universal Instantiation]
    \item[\( P(a) \lor Q(a) \)] \hfill 2, Add [Addition]
    \item[\( (\forall x)[P(x) \lor Q(x)] \)] \hfill 3, UG [Universal Generalization]
\end{enumerate}
Hence, \( (\forall x)P(x) \rightarrow (\forall x)[P(x)\lor Q(x)] \) is a valid argument.

\subsection*{Exercise 11}
Prove that the argument \( (\forall x)P(x)\land (\exists x)Q(x)\rightarrow (\exists x)[P(x)\land Q(x)] \) is valid.

\subsubsection*{Solution}
\begin{enumerate}[left=0em,labelwidth=4em,align=left]
    \item[\( (\forall x)P(x)\land (\exists x)Q(x) \)] \hfill Hyp
    \item[\( P(a) \)] \hfill 1, UI [Universal Instantiation]
    \item[\( (\exists x)Q(x) \)] \hfill Hyp
    \item[\( Q(b) \)] \hfill 1, EI [Existential Instantiation]
    \item[\( P(a) \land Q(b) \)] \hfill 2, 4, Con [Conjunction]
    \item[\( (\exists x)[P(x) \land Q(x)] \)] \hfill 5, EG [Existential Generalization]
\end{enumerate}
\subsection*{Exercise 12}
Prove that the argument \( (\exists x)(\exists y)P(x,y)\rightarrow (\exists y)(\exists x)P(x,y) \) is valid.

\subsubsection*{Solution}
\begin{enumerate}[left=0em,labelwidth=4em,align=left]
    \item[\( (\exists x)(\exists y)P(x,y) \)] \hfill Hyp
    \item[\( P(a, b) \)] \hfill 1, EI [Existential Instantiation]
    \item[\( (\exists y)P(a, y) \)] \hfill 2, EG [Existential Generalization]
    \item[\( (\exists x)(\exists y)P(x, y) \)] \hfill 3, EG [Existential Generalization]
\end{enumerate}




\newpage
\section*{Exercises 1.4.15-27}
Either prove that the wff is a valid argument or give an interpretation in which it is false.
\subsection*{Problem 15}
\textbf{Problem Statement}: Prove or disprove \( (\exists x)[A(x)\land B(x)]\rightarrow (\exists x)A(x)\land (\exists x)B(x) \)

\subsubsection*{Solution}
\begin{enumerate}[left=0em,labelwidth=4em,align=left]
    \item[\( (\exists x)[A(x)\land B(x)] \)] \hfill Hypotheses
    \item[\( A(a) \land B(a) \)] \hfill 1, Existential Instantiation
    \item[\( A(a) \)] \hfill 2, Simplification
    \item[\( B(a) \)] \hfill 2, Simplification
    \item[\( (\exists x) A(x) \)] \hfill 3, Existential Generalization
    \item[\( (\exists x) B(x) \)] \hfill 4, Existential Generalization
    \item[\( (\exists x) A(x) \land (\exists x) B(x) \)] \hfill 5, 6, Conjunction
    \item[\( (\exists x)[A(x)\land B(x)]\rightarrow (\exists x)A(x)\land (\exists x)B(x) \)] \hfill 1-7, Implication
\end{enumerate}
\newpage
\subsection*{Problem 16}
\textbf{Problem Statement}: Prove or disprove \( (\exists x)[R(x)\lor S(x)]\rightarrow (\exists x)R(x)\land (\exists x)S(x) \)

\subsubsection*{Solution}
\begin{enumerate}[left=0em,labelwidth=4em,align=left]
    \item[\( (\exists x)[R(x) \lor S(x)] \)] \hfill Hypotheses
    \item[\( R(a) \lor S(a) \)] \hfill 1, Existential Instantiation
    \item[\( (\exists x)R(x) \lor (\exists x)S(x) \)] \hfill 2, Existential Generalization

    % Split the proof into cases to show the issue
    \item[\( (\exists x)R(x) \)] \hfill 3, Simplification (Case 1)
    \item[\( (\exists x)S(x) \)] \hfill 3, Simplification (Case 2)
    
    % Implication
    \item[\( (\exists x)[R(x) \lor S(x)] \rightarrow (\exists x)R(x) \lor (\exists x)S(x) \)] \hfill 1-3, Implication
    
    % Here, we note the failure of proving the original statement
    \item[\textbf{Note:}] \hfill We cannot prove \( (\exists x)[R(x)\lor S(x)]\rightarrow (\exists x)R(x)\land (\exists x)S(x) \)
\end{enumerate}


\newpage
\subsection*{Problem 18}
\textbf{Problem Statement}: Prove or disprove \( (\forall x)(P(x))' \rightarrow (\forall x)(P(x)\rightarrow Q(x)) \)

\subsubsection*{Solution}
\begin{enumerate}[left=0em,labelwidth=4em,align=left]
    \item[\( (\forall x)\lnot P(x) \)] \hfill Hypotheses
    \item[\( \lnot P(a) \)] \hfill 1, Universal Instantiation
    \item[\( P(a) \rightarrow Q(a) \)] \hfill Implication (constructive dilemma)
    \item[\( \lnot P(a) \lor Q(a) \)] \hfill 3, Implication
    \item[\( Q(a) \)] \hfill 2, 4, Modus Ponens
    \item[\( (\forall x) Q(x) \)] \hfill 5, Universal Generalization
    \item[\( P(a) \rightarrow Q(a) \)] \hfill 5, Implication
    \item[\( (\forall x)(P(x) \rightarrow Q(x)) \)] \hfill 7, Universal Generalization
    \item[\( (\forall x)\lnot P(x) \rightarrow (\forall x)(P(x) \rightarrow Q(x)) \)] \hfill 1-8, Implication
\end{enumerate}

\newpage
\subsection*{Problem 19}
\textbf{Problem Statement}: Prove or disprove \( [(\forall x)P(x) \rightarrow (\forall x)Q(x)]\rightarrow (\forall x)[P(x)\rightarrow Q(x)] \)

\subsubsection*{Solution}
\begin{enumerate}[left=0em,labelwidth=4em,align=left]
    \item[\( (\forall x)P(x) \rightarrow (\forall x)Q(x) \)] \hfill Hypotheses
    \item[\( P(a) \)] \hfill Assume for conditional proof
    \item[\( (\forall x)P(x) \)] \hfill 2, Universal Generalization
    \item[\( (\forall x)Q(x) \)] \hfill 1, 3, Modus Ponens
    \item[\( Q(a) \)] \hfill 4, Universal Instantiation
    \item[\( P(a) \rightarrow Q(a) \)] \hfill 2, 5, Implication
    \item[\( (\forall x)[P(x) \rightarrow Q(x)] \)] \hfill 6, Universal Generalization
    \item[\( [(\forall x)P(x) \rightarrow (\forall x)Q(x)]\rightarrow (\forall x)[P(x)\rightarrow Q(x)] \)] \hfill 1-7, Implication
\end{enumerate}

\newpage
\subsection*{Problem 25}
\textbf{Problem Statement}: Prove or disprove \( (\forall x)(P(x)\lor Q(x))\land (\exists x)Q(x)\rightarrow (\exists x)P(x) \)

\subsubsection*{Solution}
\begin{enumerate}[left=0em,labelwidth=4em,align=left]
    \item[\( \text{Let } P(x): x \text{ is even}, Q(x): x \text{ is odd} \)] \hfill Counterexample Definitions
    \item[\( \forall x (P(x) \lor Q(x)) \)] \hfill Every integer is either even or odd
    \item[\( \exists x Q(x) \)] \hfill There exists an odd integer (e.g., 1)
    \item[\( (\forall x)(P(x)\lor Q(x))\land (\exists x)Q(x) \)] \hfill 2, 3, Conjunction
    \item[\( \nexists x P(x) \)] \hfill No integer exists that is both even and odd
    \item[\( (\forall x)(P(x)\lor Q(x))\land (\exists x)Q(x)\nrightarrow (\exists x)P(x) \)] \hfill 4, 5, Counterexample shown
\end{enumerate}


\newpage
\subsection*{Problem 26}
\textbf{Problem Statement}: Prove or disprove \( (\exists x)[P(x)\rightarrow Q(x)]\land (\forall y)[Q(y)\rightarrow R(y)]\land (\forall x)P(x)\rightarrow (\exists x)R(x) \)

\subsubsection*{Solution}
\begin{enumerate}[left=0em,labelwidth=4em,align=left]
    \item[\( (\exists x)[P(x)\rightarrow Q(x)]\land (\forall y)[Q(y)\rightarrow R(y)]\land (\forall x)P(x) \)] \hfill Hypotheses
    \item[\( (\exists x)[P(x)\rightarrow Q(x)] \)] \hfill 1, Simplification
    \item[\( (\forall y)[Q(y)\rightarrow R(y)] \)] \hfill 1, Simplification
    \item[\( (\forall x)P(x) \)] \hfill 1, Simplification
    \item[\( P(a)\rightarrow Q(a) \)] \hfill 2, Existential Instantiation
    \item[\( P(a) \)] \hfill 4, Universal Instantiation
    \item[\( Q(a) \)] \hfill 5, 6, Modus Ponens
    \item[\( Q(a)\rightarrow R(a) \)] \hfill 3, Universal Instantiation
    \item[\( R(a) \)] \hfill 7, 8, Modus Ponens
    \item[\( (\exists x)R(x) \)] \hfill 9, Existential Generalization
    \item[\( (\exists x)[P(x)\rightarrow Q(x)]\land (\forall y)[Q(y)\rightarrow R(y)]\land (\forall x)P(x)\rightarrow (\exists x)R(x) \)] \hfill 1-10, Implication
\end{enumerate}


\newpage
\section*{Problem 1.4.34}
\textbf{Problem Statement}:
Every computer science student works harder than somebody, and everyone who works harder than any other person gets less sleep than that person. Maria is a computer science student, therefore, Maria gets less sleep than someone else.

\subsubsection*{Solution}
\begin{enumerate}[left=0em,labelwidth=4em,align=left]
    \item[\( \forall x (C(x) \rightarrow \exists y W(x, y)) \)] \hfill Hypotheses
    \item[\( \forall x \forall y (W(x, y) \rightarrow S(x, y)) \)] \hfill Hypotheses
    \item[\( C(m) \)] \hfill Maria is a computer science student
    \item[\( C(m) \rightarrow \exists y W(m, y) \)] \hfill 1, Universal Instantiation
    \item[\( \exists y W(m, y) \)] \hfill 3, 4, Modus Ponens
    \item[\( W(m, a) \)] \hfill 5, Existential Instantiation
    \item[\( \forall y (W(m, y) \rightarrow S(m, y)) \)] \hfill 2, Universal Instantiation
    \item[\( W(m, a) \rightarrow S(m, a) \)] \hfill 7, Universal Instantiation
    \item[\( S(m, a) \)] \hfill 6, 8, Modus Ponens
    \item[\( \exists x S(m, x) \)] \hfill 9, Existential Generalization
\end{enumerate}

The conclusion, Maria gets less sleep than someone else (\( \exists x S(m, x) \)), follows from the hypotheses, so the argument is valid.


\end{document}
